% All outcome values and confidence intervals are synthetic placeholders.
% The controlled-pair benchmark and SAE audit are proposed, not completed.
% Insert before the manuscript's Robustness section.

\subsection{Groundedness Results}
\label{sec:groundedness-results}

\noindent\textit{Draft note: all results and confidence intervals below are
synthetic placeholders. The validation experiments have not been completed.}

We evaluate the Qwen2.5-7B layer-20 NLA on 1,500 held-out FineWeb documents,
using a separate 500 documents for calibration. We also select four evaluations
from the original paper's evaluation-awareness study: GPQA Diamond, MMMLU,
SWE-bench Verified, and Blackmail \citep{frasertaliente2026nla}. Each has public
questions or environments from which we can generate Qwen trajectories. These
are public replications of the evaluation settings, not the original Claude
transcripts. We use 50 tasks per evaluation, with MMMLU sampled across its
14 released languages, and report each evaluation separately.
Claim rates account for claim sampling and include unresolved cases. The
95\% intervals use 10,000 bootstrap resamples of documents for FineWeb and
tasks for the evaluations, retaining all sampled positions within each task.

\paragraph{Half of claims show specific support.}
On FineWeb, 31\% of claims are classified as grounded and 19\% as entangled
(Table~\ref{tab:grounding-population}). Both groups exceed the matched-decoy
baseline. They differ in whether the claim's own direction suppresses it more
strongly than other real-claim directions. Thus, the 31\% grounded rate alone
understates how often the test detects claim-specific support.

Another 25\% of claims are classified as seeded: the activation supports their
broader theme, but their self-ablation effects do not distinguish them from
plausible alternatives. This is nearly twice the confabulated fraction (13\%).
The more common failure identified by the test is therefore unsupported
specificity within a supported topic, rather than an entirely unrelated claim.

\begin{table*}[t]
\centering\footnotesize
\setlength{\tabcolsep}{4pt}
\begin{tabular}{lccccc}
\hline
Category & FineWeb & GPQA Diamond & MMMLU & SWE-bench Verified & Blackmail \\
\hline
Grounded & 31.0 [28.9, 33.2] & 29.0 [22.4, 36.5] & 27.0 [20.3, 34.6] & 22.0 [15.8, 29.2] & 23.0 [16.6, 30.4] \\
Entangled & 19.0 [17.3, 20.8] & 20.0 [14.7, 26.3] & 18.0 [12.8, 24.2] & 20.0 [14.2, 26.8] & 18.0 [12.6, 24.5] \\
Seeded & 25.0 [23.0, 27.1] & 25.0 [19.0, 31.8] & 24.0 [18.0, 31.0] & 26.0 [19.5, 33.4] & 29.0 [22.1, 36.7] \\
Confabulated & 13.0 [11.5, 14.6] & 14.0 [9.5, 19.5] & 15.0 [10.1, 20.9] & 14.0 [9.1, 20.0] & 13.0 [8.4, 18.7] \\
Unresolved & 12.0 [10.6, 13.6] & 12.0 [7.8, 17.3] & 16.0 [10.9, 22.0] & 18.0 [12.5, 24.6] & 17.0 [11.6, 23.5] \\
\hline
\end{tabular}
\caption{\textbf{Synthetic results.} Claim classifications (\%) with illustrative
95\% intervals. Each column sums to 100\%. The four evaluation columns represent
separate public Qwen replications; no pooled OOD estimate is reported. These
are groundedness rates over all sampled claims, not evaluation-awareness rates.}
\label{tab:grounding-population}
\end{table*}

The evaluations show different patterns. Specificity passes on GPQA Diamond
are close to FineWeb (49\% versus 50\%), compared with 45\% on MMMLU,
42\% on SWE-bench Verified, and 41\% on Blackmail. The latter two also have
more unresolved claims (18\% and 17\%, versus 12\% on FineWeb). Pooling the
evaluations would obscure this difference. Lower pass rates alone do not show
that the test is less accurate: they may reflect differences in the claims
generated, their grounding, or the availability of suitable decoys.

\paragraph{Decoys improve on claim deletion.}
Does the test identify supported details more reliably than reconstruction
alone? We construct 300 pairs of contexts that differ in one attribute, such
as a person's occupation, while preserving the topic. We score the same claim
under both activations and test whether each method distinguishes
source-consistent from source-inconsistent claims. Comparisons use the same 1,056 scorable cases out of 1,200
attempted claim--activation cases.

Our main baseline is the increase in reconstruction error after deleting a
claim, as evaluated in the original NLA paper \citep{frasertaliente2026nla}.
We compare this with the decoy-corrected self-ablation score:
\begin{equation}
\begin{aligned}
B_i^{\mathrm{del}} &= E_{\mathrm{dir}}(z_{-i};h)-E_{\mathrm{dir}}(z;h),\\
G_i &= R_i-Q_{0.95}(D_{i,1},\ldots,D_{i,20}),
\end{aligned}
\end{equation}
where $E_{\mathrm{dir}}$ is squared error between unit-normalized target and
reconstructed activations. Higher scores indicate stronger support.

\begin{table*}[t]
\centering\small
\setlength{\tabcolsep}{7pt}
\begin{tabular}{lccc}
\hline
Score & AUROC & False-positive rate (\%) & AUROC, SAE-covered \\
\hline
Whole-explanation reconstruction & .60 [.56, .64] & 10 & .61 [.56, .66] \\
Claim-deletion reconstruction & .72 [.68, .76] & 12 & .74 [.69, .79] \\
Activation-specific support, $A_i$ & .74 [.70, .78] & 13 & .76 [.71, .81] \\
Self-ablation, $R_i$ & .73 [.69, .77] & 24 & .75 [.70, .80] \\
Decoy-corrected self-ablation, $G_i$ & .82 [.78, .86] & 11 & .83 [.78, .88] \\
SAE support & --- & --- & .85 [.81, .89] \\
\hline
\end{tabular}
\caption{\textbf{Synthetic results.} Discrimination between source-consistent
and source-inconsistent claims. All non-SAE scores use the same 1,056 cases;
the last column uses the same 660 SAE-covered cases for every method.
False-positive thresholds target 10\% on a separate calibration set.
Brackets give 95\% cluster-bootstrap intervals.}
\label{tab:grounding-baselines}
\end{table*}

The corrected score reaches .82 AUROC, compared with .72 for claim deletion
(paired difference .10, 95\% CI [.06, .14]; Table~\ref{tab:grounding-baselines}).
Without decoy correction, self-ablation performs similarly to claim deletion
(.73 AUROC) and has a higher false-positive rate (24\% versus 12\%). Correction
reduces this rate to 11\%. The improvement therefore comes from comparing a
claim with plausible alternatives, not merely from measuring the AV's response
to an intervention.

This benchmark tests sensitivity to an explicit change in the source. It does
not establish that the AV faithfully reports latent beliefs or other information
that cannot be checked against the input.

\paragraph{SAEs distinguish themes from details.}
We next audit 200 natural claims from each category against an SAE at the same
model and activation site. Feature annotations are assessed without access to
our classifications. We record support for the exact claim and its broader theme separately. Support rates are computed only where
the dictionary contains an independently validated feature for the claim;
missing features are not counted as evidence against it.

\begin{table}[ht]
\centering\small
\setlength{\tabcolsep}{4pt}
\begin{tabular}{lrrr}
\hline
Category & Covered & Exact (\%) & Theme (\%) \\
\hline
Grounded & 150/200 & 82.0 & 92.0 \\
Entangled & 140/200 & 75.0 & 90.7 \\
Seeded & 120/200 & 25.0 & 83.3 \\
Confabulated & 90/200 & 20.0 & 35.6 \\
Unresolved & 100/200 & 41.0 & 62.0 \\
\hline
\end{tabular}
\caption{\textbf{Synthetic results.} Active SAE support among claims with
validated dictionary coverage. Exact and thematic support can overlap.}
\label{tab:grounding-sae}
\end{table}

SAEs support the exact content of 82\% of grounded claims and 75\% of entangled
claims (Table~\ref{tab:grounding-sae}). Entanglement therefore need not indicate
an unsupported claim. For seeded claims, thematic support remains high (83\%),
but exact support falls to 25\%. This is the distinction the decoy controls
are intended to capture: evidence for a topic need not support the particular
detail chosen by the AV.

SAEs are also a strong verifier on the controlled benchmark. On the 660 cases
with feature coverage, SAE support reaches .85 AUROC versus .83 for our score
(paired difference $-.02$, 95\% CI [$-.06$, $.02$]). We find no advantage over
SAEs on this subset. The proposed test covers additional claims, although that
broader coverage does not itself establish their reliability. Finally, 20\%
of covered claims classified as confabulated have exact SAE support. These
disagreements show why a negative result from the test should not be treated
as proof that the activation lacks the claimed information.
